\clearpage
\fancyhead[L]{\small FLORA.}
\SourceLeaf{481}
\begin{center}
\rule{0.55\linewidth}{0.4pt}\\[2pt]
Astronomische Nachrichten. Band XXVI. Nr. 622. Seite 352. 1848 Januar 10.\\[2pt]
\rule{0.55\linewidth}{0.4pt}
\end{center}

Von Herrn Geheimen Hofrath \textsc{Gauss} habe ich folgende Beobachtungen der \emph{Flora} am Göttinger Meridiankreise erhalten.

\begin{center}
\setlength{\tabcolsep}{4pt}
\begin{tabular}{l|r r r|r|r}
Dec. 10 & $11^{\rm u}$ & $11^{\rm m}$ & $4^{\rm s}0$ & $66^\circ\ 47'\ 26''2$ & $+14^\circ\ 22'\ 50''1$ \\
\hfill 11 & & $6$ & $7.9$ & $66\ 32\ 5.4$ & $14\ 25\ 42.6$ \\
\hfill 12 & & $1$ & $11.0$ & $66\ 16\ 59.9$ & $14\ 28\ 40.0$ \\
\end{tabular}
\end{center}

Der Herr Geheime Hofrath \textsc{Gauss} bemerkt, dass die AR. am 10. etwas weniger zuverlässig als die am 11. sei, und dass Flora viel heller als Iris erscheine. Die Beobachtungen konnten mit voller Beleuchtung gemacht werden, und der Planet gewährt über $\rho$ Tauri, welcher einige Minuten nördlicher ihm um $7^{\rm s}$ folgte, einen schönen Anblick.

\clearpage
\hypertarget{native-481-january4}{}
\begin{center}
\rule{0.55\linewidth}{0.4pt}\\[2pt]
Astronomische Nachrichten. Band XXVI. Nr. 623. Seite 367. 1848 Januar 17.\\[2pt]
\rule{0.55\linewidth}{0.4pt}
\end{center}
\begin{center}[\textsc{Gauss} an \textsc{Schumacher}.]\quad Göttingen 1848 Januar 4.\end{center}

Nach langer Unterbrechung hatte sich gestern der Himmel, obwohl nur auf kurze Zeit aufgeklärt, und ich bin dadurch in den Stand gesetzt worden, wieder eine Meridianbeobachtung der \emph{Flora} zu machen, die ich Ihnen hier mittheile.

\begin{center}
Beobachtung der Flora.\\[3pt]
\setlength{\tabcolsep}{4pt}
\begin{tabular}{l|r|r|r}
1848 & M. Z. in Göttingen & Ger. Aufsteig. & Abweichung \\\hline
Januar 3 & $9^{\rm u}\ 20^{\rm m}\ 4^{\rm s}2$ & $62^\circ\ 37'\ 13''8$ & $+15^\circ\ 59'\ 24''6$ \\
\end{tabular}
\end{center}

Prof. \textsc{Goldschmidts} Elemente weichen davon ab (nach seiner Rechnung)

\begin{center}
$+18''2$ in gerader Aufsteigung.\qquad $+2''4$ in der Declination.
\end{center}

\clearpage
\hypertarget{native-481-january10}{}
\begin{center}
\rule{0.55\linewidth}{0.4pt}\\[2pt]
Astronomische Nachrichten. Band XXVII. Nr. 625. Seite 13..16. 1848 Januar 29.\\[2pt]
\rule{0.55\linewidth}{0.4pt}
\end{center}
\begin{center}[\textsc{Gauss} an \textsc{Schumacher}.]\quad Göttingen 1848 Januar 10.\end{center}

Flora habe ich am 5ten im Meridian beobachtet.

\begin{center}
\setlength{\tabcolsep}{4pt}
\begin{tabular}{l|r|r|r}
1848 Januar 5 & $9^{\rm u}\ 11^{\rm m}\ 44^{\rm s}8$ & $62^\circ\ 30'\ 18''1$ & $+16^\circ\ 9'\ 45''2$ \\
\end{tabular}
\end{center}

Die Abweichung der \textsc{Goldschmidt}schen Elemente ist $+16''4$, $+3''5$. Am 3ten fand ich sie zwar bedeutend schwächer als im December v. J., aber doch noch hell genug, dass sie wohl fast volle Beleuchtung ertragen haben würde. Am 5ten musste ich die Beleuchtung sehr dämpfen. In der That vertrug der Planet kaum nur geringe Feldbeleuchtung, was ich dem Zustande der Luft zuschrieb, da auch die Fixsterne sehr unreine und unruhige Bilder gaben. Jedenfalls steht die Beobachtung vom 5ten der vom 3ten an Genauigkeit nach.
